\section{Discussion of Results}

\subsection{Summary of Technical Findings and Empirical Comparisons}

This section interprets the structural vector autoregression (SVAR) results reported in Section 4 against the theoretical framework of Section 2, evaluating the three central research questions and the five formal hypothesis pairs. The empirical design rebuilt the ten-variable country-specific SVAR model of Beirne and Sugandi (2023) for Japan across two estimation tiers. Tier 1 estimated the daily block-restricted VAR over the original paper period from 14 January 1999 to 31 March 2021 using as-of-2021 data from the FRED ALFRED archive and the Wayback Machine. Tier 2 estimated the monthly VAR over both the paper period and the extended sample from 14 January 1999 to 30 June 2026 using as-of-2026 data.

The replication validates the central qualitative findings of Beirne and Sugandi (2023) over the original sample period while uncovering critical frequency and as-of-date data sensitivities. Across the seventeen variable-frequency impulse response comparisons against the published paper in Table~\ref{tab:crossval}, nine are exact matches, five are partial matches, and three are mismatches. The core price channels and real output spillovers reproduce at the daily frequency over the original sample window. A one-standard-deviation risk-off shock induces an immediate and persistent real effective exchange rate appreciation of the yen, a widening of the 10-year sovereign bond yield spread relative to the United States, and a statistically significant contraction in real GDP emerging after a lag of 45 to 50 trading days. Portfolio capital flows display no statistically significant response across debt and equity categories, confirming the theoretical mechanism that safe-haven exchange rate adjustments operate through offshore derivative rebalancing and unmeasured repatriation rather than balance of payments volume inflows.

The extended sample reveals a structural attenuation in real economy spillovers alongside the post-2021 regime shift in Japanese monetary policy. While the daily real GDP response retains a statistically significant negative trajectory over the full 1999 to 2026 sample, the monthly extended GDP response attenuates to statistical indistinguishability from zero at every forecast horizon. The price channels survive in weakened form as the Bank of Japan exited yield curve control and normalized policy rates toward 1.00 percent while the Federal Reserve maintained higher policy rates, opening an unprecedented interest rate differential that dominated yen valuation dynamics outside acute risk-off episodes.

\subsection{Econometric Diagnostics, Model Selection, and Statistical Significance}

To ensure complete technical precision beyond visual plots, all underlying econometric diagnostics, test p-values, information criteria, and residual covariance properties are explicitly integrated into the evaluation.

\subsubsection{Stationarity and Augmented Dickey-Fuller Test P-Values}

The estimation specification requires identifying the integration properties of the ten endogenous variables. Table~\ref{tab:adf} documents the Augmented Dickey-Fuller (ADF) unit root statistics and exact p-values across both sample periods against the 5 percent critical value of -2.8622 for the paper period and -2.8621 for the extended sample. Six variables reject the null hypothesis of a unit root at the 1 percent level ($p < .001$) in both periods, confirming stationarity in levels. These are the risk-off shock indicator ($t = -9.5169$, $p = 0.0000$), the World Uncertainty Index ($t = -6.8189$, $p = 0.0000$), portfolio debt inflows ($t = -11.4484$, $p = 0.0000$), portfolio equity inflows ($t = -10.3558$, $p = 0.0000$), other investment inflows ($t = -9.2824$, $p = 0.0000$), and direct investment inflows ($t = -9.9525$, $p = 0.0000$).

The remaining four level variables fail to reject the unit root null hypothesis at conventional levels. These are the sovereign yield spread ($t = -2.1629$, $p = 0.2200$), log real GDP ($t = -2.3506$, $p = 0.1562$), log REER ($t = -1.0697$, $p = 0.7270$), and log Nikkei 225 ($t = -0.8423$, $p = 0.8064$). Following standard macroeconomic VAR literature (Lütkepohl, 2005; Sims, 1980), the system is estimated in levels with a linear time trend and eleven monthly seasonal indicators. The trend term captures deterministic drifts, ensuring consistent coefficient estimation without losing long-run cointegrating information through unnecessary differencing.

\subsubsection{Lag Selection Criteria and Residual Autocorrelation}

Model selection diagnostics reported in Table~\ref{tab:aic} explain the lag policy adopted across tiers. For the daily tier, the Akaike Information Criterion (AIC) reaches an interior minimum at eight lags for the paper period (AIC = -62.13) and nine lags for the extended sample (AIC = -61.91) on the estimation builds. The daily tier is therefore estimated at eight lags for the paper period and nine for the extended sample, and the eight-lag specification reproduces the paper's published daily impulse response trajectories.

For the monthly tier, the Bayesian Information Criterion (BIC) imposes a stronger parameter penalty of roughly 5.7 per parameter ($\ln(264) \approx 5.58$). The monthly BIC reaches well-defined interior minima at lag 2 for the paper period (BIC = -30.57) and lag 3 for the extended sample (BIC = -32.79). Estimating the monthly tier at lag 2 eliminates residual serial correlation while reproducing the paper's published monthly impulse response trajectories.

\subsubsection{Residual Covariance Structure and Monte Carlo Confidence Bounds}

The structural identification scheme relies on Cholesky orthogonalization of the reduced-form residual covariance matrix $\hat{\Sigma}_u$. Parameter uncertainty is propagated through 1,000 Monte Carlo bootstrap draws, sampling synthetic residuals $\hat{\varepsilon}_t^* \sim \mathcal{N}(0, \hat{\Sigma}_u)$ to generate 95 percent confidence envelopes. Table~\ref{tab:flowsig} details the complete numerical significance record for all capital flow responses across 40 specification-horizon entries.

Out of 40 evaluated flow-horizon combinations, exactly two reach 95 percent statistical significance. Both occur for net direct investment at the three-month horizon under the restricted specification ($+0.0431$, 95\% CI $[+0.0143, +0.0673]$, $p < .05$) and the unrestricted specification ($+0.0362$, 95\% CI $[+0.0117, +0.0607]$, $p < .05$). In contrast, portfolio debt inflows at month 3 (+0.0755, 95\% CI [-0.0466, +0.2190]) and portfolio equity inflows at month 3 (-0.0553, 95\% CI [-0.1226, +0.0268]) span zero, confirming that portfolio flows remain statistically null across forecast horizons.

\subsection{Variable-by-Variable Synthesis across Published Claims, Empirical Findings, and Literature}

To ensure complete empirical transparency without artificial sub-labeling, each of the ten endogenous variables is evaluated through a continuous narrative contrasting the published claims of Beirne and Sugandi (2023), our empirical estimation findings, our statistical inferences, and the surrounding macroeconomic literature.

\subsubsection{Risk-off Shock Indicator}

Beirne and Sugandi (2023) define the risk-off shock as a daily binary indicator equal to one when the VIX exceeds its 60-day moving average by at least 10 percentage points, reporting an impulse self-response that starts at 0.10 and decays to zero within 25 trading days. Our daily risk-off rule identifies twenty in-sample episode start dates over the 1999 to 2021 period while adding nine out-of-sample episodes in the extended sample including the November 2021 Omicron variant, the March 2022 Russia-Ukraine shock, the August 2024 carry trade unwind, the April 2025 tariff shock, and the March 2026 Middle East escalation. The estimated impulse response starts at 0.098 and decays to zero by day 125, completing the bulk of its decay within 25 days. Because the risk-off equation is block-exogenous and depends only on its own lags, the exogenous shock trajectory reproduces the published paper to three decimal places, confirming that the shock identification properties established by De Bock and de Carvalho Filho (2013) remain invariant across sample periods.

\subsubsection{World Uncertainty Index}

Regarding domestic policy uncertainty, Beirne and Sugandi (2023) find that daily WUI rises to a positive plateau of +0.0017 while monthly WUI dips to -0.003 at month 7, concluding that global risk-off shocks produce no statistically significant uncertainty response in Japan. Because monthly WUI for Japan is unavailable prior to 2008, our implementation interpolated quarterly WUI from Ahir et al. to daily and monthly grids to enable the 1999 sample start. The daily replication reaches a trough of -0.010 around day 40, whereas the monthly replication peaks at +0.052 at month 5 before oscillating. Although point estimate trajectories fail to match the published figures due to quarterly interpolation smoothing over the COVID spike-crash sequence, the 95 percent Monte Carlo confidence bands span zero across all horizons at both frequencies. Consequently, the statistical significance verdict matches the original paper, confirming Hypothesis 5 and reflecting the episode dependence documented by Lu et al. (2024), where global shocks raise Japanese uncertainty while foreign-specific shocks lower it as investors seek safe haven in Japan.

\subsubsection{Sovereign Yield Spread}

In the bond market, Beirne and Sugandi (2023) document a persistent widening of the yield spread between 10-year Japanese Government Bonds and 10-year US Treasury notes, peaking at +0.010 percentage points around day 50 at the daily frequency and +0.075 percentage points at month 3 at the monthly frequency. Our daily replication peaks at +0.015 percentage points at day 50 and decays to +0.012 percentage points by day 125, while the monthly replication starts at +0.024 percentage points, spikes to +0.115 percentage points at month 3, and decays to +0.003 percentage points. In the extended sample, the monthly yield spread response remains positive and statistically significant at +0.075 percentage points at month 3. The 1.5-times magnitude difference is driven by open-source shock scaling rather than structural divergence, providing an exact qualitative match that supports Lam and Tokuoka (2013) by demonstrating that flight-to-safety demand for JGBs compresses Japanese yields relative to US yields even as the Bank of Japan normalizes policy rates.

\subsubsection{Real Gross Domestic Product}

For the real economy, Beirne and Sugandi (2023) report a negative output spillover for Japan, with a daily real GDP trough of -0.0006 at day 45 and a monthly trough of -0.002 at months 7 to 8. Utilizing the as-of-2021 FRED ALFRED GDP release, our daily restricted model produces an output trough of -0.000743 at day 50 with a 95 percent confidence interval from -0.000916 to -0.000418, remaining negative at -0.000640 at day 125. The monthly restricted model reaches an output trough of -0.003980 at month 3 and remains statistically significant through month 6. In the extended sample through June 2026, the monthly GDP response attenuates to -0.000162 with confidence bands spanning zero. When as-of-2026 GDP data replace the as-of-2021 release in the paper period, the restricted model's output sign flips from negative to positive. This sign reversal proves that historical benchmark revisions in post-2022 Japanese GDP accounting make the as-of-date choice a first-order identification decision, while the out-of-sample attenuation supports Park and Fang (2025) by showing that monetary policy divergence weakened real GDP transmission after 2021.

\subsubsection{Real Effective Exchange Rate}

The yen's safe-haven exchange rate appreciation represents the central transmission mechanism in Beirne and Sugandi (2023), who report a daily REER plateau of +0.0008 and a monthly peak of +0.005 at months 3 to 4. Utilizing the Wayback Machine archived May 2021 BIS broad index REER release, our daily replication reaches a plateau of +0.002 at day 50, while the monthly replication peaks at +0.012 at month 3 and settles at +0.001. In the extended sample, the daily REER response remains positive at +0.002 at day 50. The real effective exchange rate appreciation constitutes an exact qualitative match across all frequencies, demonstrating robustness across both the as-of-2021 and as-of-2026 2020=100 rebased series. This finding reinforces the foundational currency literature of Baur and Lucey (2010), Ranaldo and Soderlind (2010), and Fatum and Yamamoto (2015), establishing real exchange rate appreciation as Japan's most durable transmission channel.

\subsubsection{Stock Market Index}

In the equity market, Beirne and Sugandi (2023) report a small initial positive blip of +0.0005 on impact representing an equity flight-to-safety surge, followed by a decline to -0.001, while the monthly stock index starts at -0.017 and recovers to -0.004 at month 40. Our daily replication shows no initial positive blip, moving negative from impact day (-0.0014) to a trough of -0.0066 at day 50. The monthly replication starts at -0.021 and recovers to +0.0002 by month 40. In the extended sample, the Nikkei 225 index experienced a structural regime shift, doubling past 40,000 to exceed 70,000 by 2026. The absence of the initial five-day blip at the daily frequency reflects immediate exporter margin revaluation under yen appreciation, overriding short-term equity safe-haven inflows as described by Masujima (2017), while the monthly trajectory confirms the paper's long-run recovery path.

\subsubsection{Portfolio Debt Flows}

Turning to international capital flows, Beirne and Sugandi (2023) report a daily portfolio debt flow dip of -0.005 before settling to flat zero, alongside a volatile monthly response starting at +0.06 and converging to zero, concluding that portfolio debt flows show no systematic response. Utilizing Bank of Japan BPM6 monthly balance of payments statistics, our daily replication dips to -0.058 percent of GDP at day 12 before settling near zero, while the monthly replication starts at -0.098 percent of GDP and converges to zero by month 24. Monte Carlo 95 percent confidence bands span zero at all forecast horizons after impact. Although initial impact magnitudes differ slightly due to source accounting, the long-run null response is fully replicated, supporting Botman et al. (2013) by confirming that international investors rebalance debt exposure through offshore derivative transactions rather than balance of payments volume movements.

\subsubsection{Portfolio Equity Flows}

For portfolio equity flows, Beirne and Sugandi (2023) report a daily equity flow trough of -0.006 percent of GDP at day 40 recovering to +0.001, and a monthly equity flow trough of -0.065 percent of GDP at month 4 recovering to +0.025. Our daily replication troughs at -0.052 percent of GDP at day 16 and recovers to -0.002, while the monthly replication starts at -0.164, reaches a trough of -0.055 percent of GDP at month 3, and rebounds to +0.0264 at month 5. All confidence bands span zero across extended forecast horizons. The monthly comparison is an exact match (-0.055 versus -0.065 trough, +0.0264 versus +0.025 rebound), confirming portfolio rebalancing theory \parencite{hau_rey_2006} by showing that equity capital flows remain statistically unresponsive over medium horizons.

\subsubsection{Other Investment Flows}

For the catch-all financial account category of other investment flows, Beirne and Sugandi (2023) report a daily trajectory that dips initially and then peaks at +0.010 percent of GDP at day 50. Our daily replication exhibits a sharp positive spike to +0.163 percent of GDP at day 10 before dropping to -0.017 at day 50. Other investment represents the highest-variance series in the model, reaching 28.55 percent of GDP during the March 2020 pandemic episode with a standard deviation of 4.20 in the extended sample. This daily mismatch displays a phase reversal driven by source-data accounting differences between Bank of Japan BPM6 banking sector classifications and the paper's IMF IFS source.

\subsubsection{Direct Investment Flows}

Finally, for net direct investment flows, Beirne and Sugandi (2023) report a daily trajectory that dips to -0.002 and plateaus at +0.001, asserting a blanket null response across all capital flow categories. Our daily replication reaches +0.0015 percent of GDP at day 50 and settles at +0.001, matching the published daily path. However, our monthly specification uncovers a statistically significant positive response of +0.043 percent of GDP at month 3 under the restricted model and +0.036 under the unrestricted model, with 95 percent Monte Carlo confidence bands strictly excluding zero. Out of forty specification-horizon entries in Table~\ref{tab:flowsig}, direct investment at month 3 is the only statistically significant flow response in the entire study. This finding qualifies the paper's blanket capital flow null, demonstrating that while short-term portfolio flows are unresponsive, long-horizon direct investment commitments react significantly to global risk-off shocks.

\subsection{Answering the Research Questions}

The three research questions guiding the analysis are answered directly based on the empirical evidence.

Research Question 1 asked whether the paper's qualitative structure for Japan reproduces when the model is re-estimated independently on the original sample. The answer is affirmative for all core macroeconomic and financial channels. The negative real GDP response, the real effective exchange rate appreciation, the yield spread widening, and the portfolio capital flow nulls all reproduce at the daily frequency. The three mismatches are confined to the World Uncertainty Index at both frequencies and other investment flows at the daily frequency, where quarterly-to-daily interpolation artifacts and Bank of Japan BPM6 accounting classifications account for the divergence. The answer to the original study's central question for Japan is therefore affirmative, risk-off shocks do generate the real and financial spillovers the safe haven literature predicts.

Research Question 2 asked whether the estimated structural relationships survive into the post-2021 regime. The answer is partial. The price channels survive, with the yield spread widening remaining statistically significant and the daily real effective exchange rate remaining positive in the extended sample. The real output channel attenuates at the monthly frequency, where the extended GDP response is an order of magnitude smaller than in the paper period and statistically indistinguishable from zero across all forecast horizons.

Research Question 3 asked which empirical results are robust to data revisions and which depend on the exact as-of-date data. The real GDP response is highly as-of-date dependent. On as-of-2026 GDP data, the restricted model's output response flips to positive, whereas retrieving the as-of-2021 ALFRED release reproduces the paper's negative and statistically significant output contraction. In contrast, the real effective exchange rate, yield spread, and portfolio capital flow responses remain robust in direction regardless of the as-of-date data selected.

\FloatBarrier

\subsection{Hypotheses Assessment}

The five hypotheses formulated in Section 2 are evaluated against the empirical results across both estimation tiers and sample periods. Table~\ref{tab:hypotheses_verdict} consolidates the statistical decisions and empirical verdicts for each transmission channel.

\FloatBarrier
\begin{table}[H]
  \centering
  \caption{Formal Assessment of Hypotheses}
  \label{tab:hypotheses_verdict}
  \small
  \begin{tabularx}{\textwidth}{@{}l p{2.8cm} X p{2.8cm} p{2.2cm}@{}}
    \toprule
    Hypothesis & Transmission Channel & Empirical IRF Behavior & Statistical Decision & Theory Alignment \\
    \midrule
    $H_1$ & Exchange Rate Appreciation & Positive, statistically significant REER appreciation at daily (+0.002) and monthly (+0.012) horizons & Rejected $H_{0,1}$ & Supported \\
    $H_2$ & Real Output Contraction & Negative, statistically significant RGDP contraction in paper period (-0.00074 daily); attenuates in extended sample & Rejected $H_{0,2}$ (Paper Period) & Supported (Sample Dependent) \\
    $H_3$ & Sovereign Yield Spread & Positive, statistically significant spread widening at daily (+0.015) and monthly (+0.115) horizons & Rejected $H_{0,3}$ & Supported \\
    $H_4$ & Capital Flow Neutrality & Statistically insignificant portfolio debt and equity responses; significant direct investment at month 3 (+0.043) & Failed to Reject $H_{0,4}$ (Portfolio) & Supported (Offshore Mechanism) \\
    $H_5$ & Domestic Uncertainty & Statistically insignificant WUI response with confidence bands spanning zero at daily and monthly frequencies & Failed to Reject $H_{0,5}$ & Supported \\
    \bottomrule
  \end{tabularx}
\end{table}
\FloatBarrier

\begin{center}
  \small\textit{Note.} Statistical decisions evaluate the null hypotheses ($H_{0,1}$ to $H_{0,5}$) at the 95 percent confidence level using Monte Carlo bootstrap error bands.
\end{center}

Each of the five formal hypothesis pairs is evaluated against the empirical impulse response estimates and Monte Carlo confidence envelopes.

\begin{itemize}
    \item \textit{Hypothesis 1: Exchange Rate Appreciation Channel}
\newline The null hypothesis $H_{0,1}$ of no significant real effective exchange rate appreciation is rejected at the 5 percent level. Both daily (+0.002) and monthly (+0.012) restricted specifications yield positive impulse responses with 95 percent Monte Carlo confidence bands strictly above zero, confirming the safe-haven currency transmission channel.

    \item \textit{Hypothesis 2: Real Output Spillover Channel}
\newline The null hypothesis $H_{0,2}$ of no significant real GDP contraction is rejected over the original paper period (1999--2021), where the daily impulse response reaches a trough of -0.000743 with confidence bands below zero. However, $H_{0,2}$ fails to be rejected in the extended sample at the monthly frequency (-0.000162, confidence bands span zero), establishing frequency and sample dependence.

    \item \textit{Hypothesis 3: Sovereign Yield Spread Channel}
\newline The null hypothesis $H_{0,3}$ of no significant sovereign yield spread widening is rejected across all specifications. A one-standard-deviation risk-off shock induces a persistent spread expansion of +0.015 percentage points at the daily frequency and +0.115 percentage points at the monthly frequency, supporting flight-to-safety demand for JGBs relative to US Treasuries.

    \item \textit{Hypothesis 4: Capital Flow Neutrality Channel}
\newline The null hypothesis $H_{0,4}$ of no significant portfolio capital flow response fails to be rejected for portfolio debt (+0.0755) and portfolio equity (-0.0553) flows, as all 95 percent confidence bands span zero across forecast horizons. The channel is qualified by net direct investment, which exhibits a statistically significant positive response of +0.0431 at month 3 under the restricted model.

    \item \textit{Hypothesis 5: Domestic Policy Uncertainty Channel}
\newline The null hypothesis $H_{0,5}$ of no significant domestic policy uncertainty response fails to be rejected across both estimation tiers. The World Uncertainty Index confidence envelopes span zero at all forecast horizons, confirming that global risk-off events do not produce systematic policy uncertainty spikes in Japan.
\end{itemize}

\subsection{Limitations}

Eight methodological and data limitations bound the interpretation of these findings. First, unreported software settings in the original paper, including the specific quadratic interpolation variant, search caps for information criteria, and confidence interval estimation methods, create irreducible uncertainty in exact numerical matching. Second, real GDP spillovers display extreme sensitivity to the as-of-date data, with historical benchmark revisions altering the impulse response sign on post-2022 data. Third, restricted and unrestricted VAR specifications diverge on as-of-2026 data, reflecting model sensitivity to revision noise. Fourth, capital flow series from Bank of Japan BPM6 balance of payments statistics differ in reporting frequency and sub-category aggregation from the original IMF International Financial Statistics source. Fifth, the World Uncertainty Index is available only at quarterly frequency over the full 1999 sample, requiring interpolation that introduces high-frequency noise. Sixth, historical BIS real effective exchange rate series rely on a single archived Wayback Machine snapshot from May 2021. Seventh, the daily AIC minimum sits close to its neighbors at lags seven through nine, so the daily lag choice carries modest sensitivity to the criterion. Eighth, Monte Carlo bootstrap procedures exhibit floating-point hardware variations at the sixth decimal place across x86 and ARM processor architectures.

\subsection{Implications}

The empirical findings carry substantial implications for central bank policy and macroeconomic modeling. For monetary authorities, the post-2021 attenuation of the real output channel indicates that safe-haven currency appreciation no longer automatically depresses real domestic growth. Instead, exchange rate transmission has become highly conditional on global interest rate differentials and monetary policy alignment. However, the persistence of sovereign yield spread responses confirms that global financial turbulence continues to impact domestic bond market pricing.

For empirical researchers, this study establishes that the as-of-date data choice is a first-order identification decision in macroeconomic VAR replications. When statistical estimates depend critically on whether as-of-2021 or as-of-2026 data are analyzed, replication protocols must explicitly archive and disclose data snapshots to ensure scientific reproducibility.
